% =====================================================================
\clearpage
\part{The wall in two coordinates}
% =====================================================================
\section{Coordinate 1: measure and the weak-$*$ discontinuity}
\label{ch:weakstar}
% =====================================================================

\subsection{The coordinate}

The first coordinate is the one the campaign used from the beginning: the gauge-invariant
spectral measure $\nu$ of $\operatorname{tr}F^2$, its atom $f_1^2$ at the bottom of the support,
and its moments. This chapter sharpens the structural reason every tool built in that coordinate
failed: the atom is precisely the failure of weak-$*$ continuity of the boundary density, and
every finitely computable datum is weak-$*$ continuous.

We work on the space $\mathcal M$ of finite positive measures on $[0,\infty)$ with bounded mass,
equipped with the weak-$*$ topology against bounded continuous test functions. This is the
natural topology of ``finite data'': a measurement is an integral against a bounded continuous
kernel, and convergence in this topology is exactly convergence of all such integrals.

Let
\begin{equation}
\Pi(z)=\int_{[0,\infty)}\frac{d\nu(s)}{s-z},\qquad \operatorname{Im}z\ne0 ,
\label{eq:stieltjes}
\end{equation}
be the Stieltjes transform (the boundary density, up to normalization). For an atom
$w\delta_m$, one has $\Pi(z)=w/(m-z)$, hence on the boundary $z=x+i\varepsilon$,
\begin{equation}
\frac{1}{\pi}\operatorname{Im}\Pi(x+i\varepsilon)
=\frac{w}{\pi}\frac{\varepsilon}{(x-m)^2+\varepsilon^2}
\;\longrightarrow\; w\,\delta_m(x)\quad(\varepsilon\downarrow0),
\label{eq:atomjump}
\end{equation}
so the atom weight is the boundary jump. The window must be taken independently of the
regularization and the limits must be ordered:
\begin{equation}
f_1^2=\nu(\{m\})=\lim_{d\downarrow0}\ \lim_{\varepsilon\downarrow0}
\int_{m-d}^{m+d}\frac{1}{\pi}\operatorname{Im}\Pi(x+i\varepsilon)\,dx .
\label{eq:atomweight}
\end{equation}
The limit is essential: for every fixed $\varepsilon>0$, the expression on the right is an
integral of a bounded kernel, i.e.\ a finite datum. The atom appears only in the limit.
The order matters: setting the window equal to the regulator, $d=\varepsilon$, couples the two
limits and returns $f_1^2/2$, because the Poisson kernel integrates to
$(1/\pi)[\arctan(1)-\arctan(-1)]=1/2$ over $[m-\varepsilon,m+\varepsilon]$. The correct statement
is $d\downarrow0$ \emph{after} $\varepsilon\downarrow0$; for every fixed $d>0$ the inner limit is
the full atom weight.

\subsection{Four lemmas}

The following four statements are the mathematical core of the coordinate. They are elementary;
the campaign's contribution is to recognize that they classify every tool it built.

\begin{claimbox}{cT}
\claimhead{Lemma}{(finite data is weak-$*$ continuous)} \emph{Let $\varphi_1,\dots,\varphi_J\in
C_b([0,\infty))$. The map $\nu\mapsto(\int\varphi_1\,d\nu,\dots,\int\varphi_J\,d\nu)$ is
continuous in the weak-$*$ topology. In particular: (i) point evaluations $\Pi(z_j)$ at
$\operatorname{Im}z_j\ne0$ are continuous, since $|1/(s-z_j)|\le1/|\operatorname{Im}z_j|$;
(ii) Taylor coefficients of $\Pi$ at such a point are continuous; (iii) integrals against
bounded continuous kernels are continuous.}
\end{claimbox}

\begin{claimbox}{cT}
\claimhead{Lemma}{(atomless measures are weak-$*$ dense)} \emph{For every finite measure $\nu$ on
$[0,\infty)$ and every finite family $\varphi_1,\dots,\varphi_J\in C_b$ and every $\eta>0$,
there is an atomless measure $\nu'$ with $\sup_j|\int\varphi_j\,d(\nu-\nu')|<\eta$ and
$\nu'([0,\infty))=\nu([0,\infty))$.}
\begin{proofx}
Approximate each atom $\delta_{x}$ by the uniform measure on $[x-h,x+h]$; by uniform continuity
of the finitely many $\varphi_j$ on the (compact) region carrying the mass, the integrals move by
at most $\eta$ for $h$ small. The result is atomless by construction.
\end{proofx}
\end{claimbox}

\begin{claimbox}{cT}
\claimhead{Lemma}{(the atom weight is discontinuous)} \emph{The functional $\nu\mapsto\nu(\{x\})$
is not continuous at any measure with $\nu(\{x\})>0$.}
\begin{proofx}
Take the bump replacement $\nu_h$ of Lemma 2 at the point $x$. Then $\nu_h\to\nu$ weak-$*$ while
$\nu_h(\{x\})=0$ and $\nu(\{x\})=w>0$.
\end{proofx}
\end{claimbox}

\begin{claimbox}{cT}
\claimhead{Theorem}{(no continuous forcing functional)} \emph{Let $F:\mathcal M\to\mathbb R$ be
weak-$*$ continuous and suppose $F(\nu)=0$ for every atomless $\nu$. Then $F\equiv0$. In
particular, no weak-$*$ continuous functional can be zero on the free massless field and
strictly positive on a measure with a positive atom.}
\begin{proofx}
The atomless measures are weak-$*$ dense (Lemma 2), and $F$ is continuous, so $F$ vanishes on
the closure of the atomless set, which is all of $\mathcal M$.
\end{proofx}
\end{claimbox}

\subsection{The topology, precisely}
\label{sec:topology}

The four lemmas are stated for the topology generated by bounded continuous test functions. For
rigour, it is enough to work on measures supported in a fixed compact interval
$K=[0,S]$, because the atom and the low moments live in a bounded window; the extension to
$[0,\infty)$ is by tightness. On $K$ the space $\mathcal M(K)$ of finite positive measures is the
positive cone of the dual $C(K)^*$, and the \emph{weak-$*$ topology} is the topology of
convergence of $\int\varphi\,d\nu$ for every $\varphi\in C(K)$.

\begin{itemize}
\item \textbf{Compactness.} By Banach--Alaoglu, the unit ball of $C(K)^*$ is weak-$*$ compact;
  the set of positive measures with mass $\le M$ is weak-$*$ closed in that ball, hence
  weak-$*$ compact.
\item \textbf{Continuity of finite data.} For every fixed $\varphi\in C(K)$ --- in particular
  every polynomial $s^k$, $k\le K$, and every point-evaluation kernel $1/(s-z)$ with
  $\operatorname{Im}z\ne0$ --- the map $\nu\mapsto\int\varphi\,d\nu$ is weak-$*$ continuous by
  definition of the topology.
\item \textbf{The atom is not continuous.} For $x\in K$ with $\nu(\{x\})=w>0$, the bump
  replacement $\nu_h$ of Lemma 2 converges weak-$*$ to $\nu$ while $\nu_h(\{x\})=0$. So
  $\nu\mapsto\nu(\{x\})$ is not weak-$*$ continuous.
\item \textbf{The full moment sequence is injective but not a homeomorphism.} By the Hausdorff
  moment problem, a finite measure on $[0,S]$ is determined by \emph{all} its moments
  $(M_k)_{k\ge0}$. So the atom is a function of the full moment sequence. But the inverse map is
  not continuous: along the bump family, all moments converge while the atom jumps. Any
  \emph{finite} truncation is continuous (bullet 2) and therefore blind (Theorem 4). This is the
  precise meaning of the campaign's slogan ``infinite data $=$ H1'': the full sequence determines
  the measure, but its determination is discontinuous, and every finite truncation is
  continuous.
\end{itemize}

The distinction between ``weak-$*$ continuous'' and ``weak-$*$ discontinuous'' is therefore not a
technicality. It is the exact line between the data a finite computation can supply and the atom
it cannot.

\textbf{The acid test is therefore a discontinuity test.} The campaign's free-field acid test ---
a tool must vanish on the massless free field, where $f_1^2=0$, and be positive on Yang--Mills ---
is exactly the demand that the tool be a weak-$*$ discontinuous functional. Any tool built from
finitely many moments, Taylor coefficients, point evaluations, bounded-kernel integrals, or
Wilson-loop images is weak-$*$ continuous and cannot pass. This is not a heuristic; it is the
content of Theorem 4.

\figweakstar

\subsection{The falsification wave: eight attempts to escape}

Theorem 4 says a continuous functional cannot force the atom. The natural question is whether a
discontinuous one can, without being merely a reading. The falsification wave built eight
candidates, each designed to be ``neither YM-structure nor infinite data.'' All eight fail, and
the pattern of failure is uniform (\code{ONLY-EXIT-SYNTHESIS.md}).

\begin{table}[t]
\centering
\footnotesize
\begin{tabularx}{\textwidth}{@{}>{\raggedright\arraybackslash}p{2.9cm}
  >{\raggedright\arraybackslash}p{3.9cm}>{\raggedright\arraybackslash}X@{}}
\toprule
\textbf{Route (file)} & \textbf{Attempted counterexample} & \textbf{Why it fails}\\
\midrule
Hardy upgrade (\code{FALSIFY-HARDY-UPGRADE}) & push the $L^2$/singular kernel from presence to
  forcing the weight & \LT{} over any fixed window the kernel is bounded (OPE-RHO blind); the
  escape lives only at the unknown pole, and ``singular kernel $\iff$ location $+$
  $\varepsilon\to0$'' cannot be decoupled\\
Higher singularity (\code{FALSIFY-HIGHER-SINGULAR}) & $\int|G^{(k)}|^2$, $\int|G|^p$ --- a
  stronger singularity forces the weight & \LT{} the order $p$ changes the divergence power
  $\varepsilon^{-(p-1)}$, not the need for the pole location and $\varepsilon\to0$ (= A3)\\
Entropy (\code{FALSIFY-ENTROPY}) & differential/relative entropy, Fisher, R\'enyi & convex
  $\Rightarrow$ blind; density-requiring $\Rightarrow$ diverges only at $\varepsilon\to0$;
  presence only; Markov $\Rightarrow$ BE\\
OPE coefficients (\code{FALSIFY-OPE}) & force from perturbative-UV OPE data & \LT{} $\Pi$ is
  linear in $\nu$, each OPE coefficient is a moment $M_{-k}$ $\Rightarrow$ no nonlinear
  constraint; finite OPE $=$ OPE-RHO blind, full $=$ H1\\
Nonlinear Ward/anomaly (\code{FALSIFY-WARD-NONLINEAR}) & nonlinear combinations of Ward, anomaly,
  positivity, AF tail & \LT{} an atomless flat bump matches every moment with $f_1^2=0$ and
  satisfies all exact identities; any continuous function of them is blind\\
Topological/winding (\code{FALSIFY-TOPOLOGICAL}) & winding / argument principle of the Herglotz
  function & \LT{} off-axis the winding is identical (blind); it discriminates only across
  $\mathbb{R}$, where the atom is a weak-$*$ discontinuity\\
Spectral flow (\code{FALSIFY-FLOW}) & flow of a deformed transfer operator & \LT{} flow is a
  spectrum property, the atom is an overlap; independent (nonzero flow coexists with
  $f_1^2=0$); rigid on connected components\\
Field-level (\code{FALSIFY-FIELD}) & functionals of the field beyond $\nu$ (Wilson loops,
  $\langle\operatorname{tr}F^n\rangle$, instantons, gradient flow) & \LT{} (A) correlator
  functionals $=\nu$, OPE-RHO blind; (B) beyond-$\nu$ data are 1-form/topological,
  Schur-orthogonal to $\operatorname{tr}F^2$; the bridge is a dynamical theorem $=$ H1\\
\bottomrule
\end{tabularx}
\caption{The eight falsification routes. None builds a forcing functional that is neither a
reading nor an infinite-data limit.}
\label{tab:falsify}
\end{table}

\textbf{The sharpening: one exit, not two.} The field-level route closes the last hope --- that
``use the specific YM structure'' is a separate loophole. A field functional that is a function
of the $0^{++}$ correlator is a function of $\nu$, hence an integral, hence blind. A field
functional using data beyond $\nu$ (Wilson loops, topology) is 1-form or topological,
Schur-orthogonal to the $0^{++}$ scalar, and carries no kinematic information on $f_1^2$; linking
it to $f_1^2$ requires a dynamical theorem about the specific YM measure --- which is H1. Hence
``YM structure'' is not a separate exit. There is one exit, reachable by two coordinates (a
YM-specific nonlinear constraint, which is H1; or infinite data $J\to\infty$, which is H1).

\subsection{The finite-volume corollary}

An important corollary, which reconciles the numerics with the structural picture: at finite
$L$, the atom is \emph{trivially} $f_1^2>0$, because the transfer matrix is compact and the
spectrum is discrete. So a finite-lattice measurement of $f_1^2>0$ is vacuous as regards Clay;
it is a statement about a finite-dimensional operator, not about the continuum theory. The Clay
content is the \emph{$a\to0$, $L\to\infty$ survival and isolation} of the atom, which is
isolation I2 $=$ H1. This is why Chapter~\ref{ch:numerics} reports finite-volume numbers with
the caveat that they are finite-volume facts, and why the diagnostic chapter insists that a
single volume cannot decide.

\subsection{What the coordinate proves and does not prove}

\textbf{Proved.} (i) The four lemmas, hence: no bounded-kernel integral functional can force the
atom (Theorem 4). (ii) The classification of the campaign's tools: every tool whose output is a
function of finitely many bounded-kernel integrals is in the blind class; every tool that reads
zeros or poles is A3; every support/geometric tool sees only the support; the dual/LP optimum is
zero because the atomless measure is feasible. (iii) The one escape for presence: the Hardy
$L^2$ singular-kernel functional (Chapter~\ref{sec:hardy}).

\textbf{Not proved.} That \emph{no} discontinuous functional can force the atom with YM-specific
information. On the contrary, such a functional would be a proof of H1, and the campaign's
failure to build one is the open problem. The theorem is a necessary condition on the \emph{type}
of tool, not a no-go for the gap.

\textbf{The scoped claim.} The falsification wave's conclusion is: within the eight attempted
routes, every forcing functional is either a reading or an infinite-data limit; the weak-$*$
argument is general but the ``every tool is in this class'' statement is evidence from the wave,
not a universal theorem. The paper labels it accordingly.

\subsection{What a proof in this coordinate would look like}

A proof of the gap in the measure coordinate would be a functional $F$ on measures such that:
\begin{enumerate}
\item $F(\nu_{\text{free}})=0$ and $F(\nu_{\mathord{:}\varphi^2\mathord{:}})=0$ (acid-clean on
  both atomless adversary classes);
\item $F(\nu_{\text{YM}})\ge c>0$ for the Haar-carrying interacting measure at $\beta\approx2.5$,
  nonperturbatively (no expansion, no small parameter);
\item $F$ is $(L,a)$-uniform and survives the continuum limit;
\item $F$ is not a function of finitely many moments (Theorem 4), and its evaluation either
  reads a pole (A3) or uses infinite data (H1).
\end{enumerate}
The negative-association inequality $\operatorname{Cov}(G,\Phi'|0)\le0$ is the sharpest candidate
of this type: it is deterministic, it uses the specific frustrated measure, and it is absent from
every known dependence class. Building it, or proving it impossible, is the open problem. This is
the ``one exit'' in the measure coordinate.

\subsection{The read-outs, and why they are not proofs}

The campaign did build genuine read-outs of the atom: the P3 approach law
$T(\tau)=-d\log|m-\hat e_1|/d\tau$, which tends to $\Delta$ for an atom and to $1/\tau$ for a cut;
the one-sided isolation bound $\Delta\ge\kappa(\tau)$; the Re-slope statistic
$d\operatorname{Re}\tau/d\ln c_1=1/E_1$; and the finite-size law
$\Delta(L)=m_1+c/(m_1L^2)$. Chapter~\ref{ch:numerics} states them with their protocols. Their
logical status is uniform: each is a read, not a force. The approach law needs the gap to
normalize, the isolation bound needs the edge, the Re-slope needs exact complex-$\tau$ zeros
(A3), and the finite-size law needs $L\to\infty$. All four are therefore gated by H1. They are
the best available instruments for \emph{seeing} the atom, and the coordinate's theorem
(Theorem 4) explains why no instrument of this kind can \emph{force} it.
